\chapter{Conclusion}\label{ch:concl}

\section{Summary of Contributions}

This work presented a complete machine learning system for predicting the
performance of photoelectrochemical water-splitting devices from material and
experimental parameters.
The main contributions are as follows.

\begin{enumerate}
  \item \textbf{Curated and augmented PEC dataset.}
        A dataset of 676 entries was assembled from 529 published experimental
        results and extended through two complementary strategies: physics-based
        estimation of STH and H\textsubscript{2} values from measured photocurrents
        (using the STH formula and Faraday's law), and 147 literature-anchored
        synthetic rows covering 25 well-characterised material systems.
        Augmentation increased STH coverage from 9\% to 82\% and
        H\textsubscript{2} coverage from 11\% to 83\%.

  \item \textbf{Physics-informed feature engineering.}
        Twenty numerical features were engineered from eight raw inputs,
        including interaction terms motivated by the Nernst equation
        ($E_g \times \mathrm{pH}$), Butler--Volmer kinetics
        ($V_{bias} \times I_{proxy}$), and Arrhenius-like thermal activation
        ($\mathrm{pH} \times T$).
        These features encode electrochemical domain knowledge that is not
        derivable from raw inputs alone.

  \item \textbf{Ensemble ML pipeline with uncertainty quantification.}
        XGBoost and Random Forest regressors are trained independently per
        target using 5-fold cross-validation.
        The ensemble inference strategy weights predictions by CV $R^2$
        (Equation~\eqref{eq:ensemble}) and computes a confidence score from
        model quality and inter-model agreement (Equation~\eqref{eq:confidence}).

  \item \textbf{Data-driven performance classification.}
        A composite score calibrated to training-set percentile thresholds
        and modulated by the confidence factor replaces the previous
        hardcoded threshold, enabling the HIGH class to be predicted for
        genuinely top-quartile inputs.

  \item \textbf{Deployed full-stack application.}
        A Flask REST API and a React single-page web interface expose the
        trained models with real-time predictions, confidence scores,
        and performance classification, deployed on a production server
        via Gunicorn and Nginx and accessible at \url{http://178.105.34.179/}.

  \item \textbf{Measurable model improvement through augmentation.}
        Cross-validated $R^2$ improved from 0.065 to 0.454 for Photocurrent
        Density, from 0.278 to 0.533 for STH Efficiency, and from $-$0.090
        to 0.366 for H\textsubscript{2} Evolution Rate after augmentation
        and physics-informed feature expansion.
\end{enumerate}

\section{Limitations and Future Work}

Several directions remain open for future investigation.

\begin{itemize}
  \item \textbf{Expanded real dataset.}
        The most impactful improvement would be increasing the number of genuine
        experimental STH and H\textsubscript{2} measurements.
        Automated NLP-based table extraction tools (e.g., GROBID, ChemDataExtractor)
        applied to open-access literature could substantially expand the dataset
        within a single research semester.

  \item \textbf{DFT-derived features.}
        The Materials Project API~\cite{jain2013commentary} provides free access
        to computed formation energies, effective carrier masses, and dielectric
        constants for thousands of inorganic compounds.
        Adding even two or three of these as features would give the model
        material-specific information that composition labels alone cannot carry ---
        two BiVO\textsubscript{4} samples synthesised by different routes can have
        very different carrier mobilities, but both currently map to the same string
        in the feature space.
        AFLOW~\cite{curtarolo2012aflow} is an alternative source for materials with
        less coverage in the Materials Project.

  \item \textbf{Advanced models.}
        Tree ensembles work well at our current data size but the confidence score
        we compute is a proxy, not a statistically calibrated uncertainty estimate.
        Once genuine experimental data grows past roughly 2,000 rows, Bayesian
        neural networks or MC dropout ensembles become viable and would give proper
        predictive intervals rather than the agreement-based confidence measure
        used here.
        Graph neural networks operating on crystal structures are a longer-term
        option, contingent on consistent structural data being available across
        the material systems in the dataset.

  \item \textbf{Multi-target learning.}
        A joint model predicting $J_{ph}$, STH, and HER simultaneously could
        exploit the known thermodynamic correlations between targets and
        potentially reduce data requirements for the sparse HER target.

  \item \textbf{Active learning loop.}
        The confidence score could drive an active learning workflow:
        configurations with low confidence and predicted high performance
        would be recommended to experimentalists, closing the loop
        between prediction and data acquisition.
\end{itemize}

Despite current limitations, the system provides a practical and immediately
deployable tool for accelerating the pre-screening of PEC photoelectrode
candidates.
By reducing the number of synthesise-and-test iterations before identifying
high-efficiency configurations, the system contributes to the broader goal of
accelerating the development of sustainable solar hydrogen production.
